%% ===========================================================================
\FloatBarrier
\section{Particular cases}
\label{sec:limits}
%% ===========================================================================

The general model of Sec.~\ref{sec:formulation} contains four switches, and
each of them has to reproduce a known theory when it is turned off. This
section states those reductions; the numerical results that follow are read
against them, and Sec.~\ref{sec:verification} checks two of them as exact
identities rather than as trends.

\FloatBarrier
\subsection{Classical piezoelectricity: $M_{mn}\to0$ and $\mu_{kijl}\to0$}
\label{sec:limit_classical}

If the averaging window shrinks to a point, $M_{mn}\to0$, the gradient energy
in \eqref{eq:energy} vanishes. That alone does not remove the higher-order
stress: by \eqref{eq:tau} it reduces to $-\mu_{kijm}E_k$, so a model with a
window of zero size and a nonzero flexoelectric coefficient keeps a
higher-order stress, a double traction and a fourth-order character. Classical
piezoelectricity is therefore the double limit $M_{mn}\to0$ \emph{and}
$\mu_{kijl}\to0$, in which the double-traction condition of Sec.~\ref{sec:bc}
drops out of the formulation. What remains is the classical piezoelectric strip: the weak
form keeps only the elastic, piezoelectric and dielectric terms, the
displacement field needs $C^0$ continuity only, and the coupling factor is the
purely piezoelectric one. The window size is therefore the single parameter
that switches the model from a local electromechanical theory to a
strain-gradient one, and it does so without changing any other term.

\FloatBarrier
\subsection{Purely piezoelectric response: $\mu_{kijl}=0$}
\label{sec:limit_piezo}

With the flexoelectric tensor set to zero the electric displacement
\eqref{eq:D} loses its strain-gradient contribution and the model reduces to
the linear piezoelectric strip of the plane-stress reduction
\eqref{eq:Cp}, with the strain-gradient stiffness retained. On the extensional configuration of
Fig.~\ref{fig:geometry} the coupling factor is $k^2=0.5581$ and is unchanged by
the switch, because that configuration collects the piezoelectric charge only.
On the bending configuration the piezoelectric charge of the mode integrates
to zero over a face electrode, and the piezoelectric-only coupling factor of
that mode falls to $1.0\times10^{-10}$, which is the round-off floor of the
secular equation rather than a physical value. The comparison of the two
configurations therefore cannot be read as a competition between the two
effects; it is a statement about which electrode pair can collect which
polarization.

\FloatBarrier
\subsection{Purely flexoelectric response: $e_{kij}=0$}
\label{sec:limit_flexo}

Removing the piezoelectric coupling instead leaves the flexoelectric bending
response of the same strip. Over the thickness range
$\SI{10}{\micro\metre}\to\SI{1}{\micro\metre}$ the flexoelectric coupling of
that configuration then rises from $1.76\times10^{-4}$ to
$5.5\times10^{-3}$, a factor $31$. With the piezoelectric terms present the
same sweep gives a rise from $1.74\times10^{-4}$ to $4.2\times10^{-3}$, a
factor $24$. The piezoelectric terms are electrically silent on that device,
but they are not mechanically silent: they change the elastic stiffness of the
mode that carries the gradient, and the difference between $24$ and $31$ is
the measure of that change. Both are reported, and the paper quotes the factor
$24$ because it is the response of the real crystal.

\FloatBarrier
\subsection{Classical Kelvin--Voigt and undamped limits: $\alpha=1$ and
$\tau\to0$}
\label{sec:limit_memory}

For $\alpha=1$ the fractional element \eqref{eq:kvfreq} becomes the classical
Kelvin--Voigt element, whose exact damping ratio is $\omega\tau/2$. The
implementation reproduces it: at $\omega\tau=0.3$ the ring-down of
Sec.~\ref{sec:memorystudy} returns $\zeta_{\rm eff}=0.15003$ against the exact
$0.15000$, a difference of $0.02\,\%$. For $\tau\to0$ the memory term of
\eqref{eq:kv} vanishes, the time-discrete problem \eqref{eq:step} becomes the
undamped one, and the eigenvalues of the pencil are those of the short-circuit
matrix. Between these two limits the model is genuinely fractional: the
damping ratio is no longer proportional to $\omega\tau$, and the memory matrix
that the derivative multiplies changes $\zeta_{\rm eff}$ by a factor $0.450$
at $\alpha=1$ (Sec.~\ref{sec:memorystudy}).

\FloatBarrier
\subsection{Recovery of the one-scale gradient theory}
\label{sec:limit_fem1}

Finally, the one-scale gradient theories are recovered as special choices of
$M$. For an ellipsoidal window $M_{mn}=L_{mn}/5$, so
$\tfrac14M_{mn}=\tfrac1{20}L_{mn}$ and the tensorial strain-gradient energy of
the one-scale theory~\cite{fem1} is reproduced exactly, with the cuboid,
ellipse and ellipsoid windows giving the scalar coefficients that
\ref{app:Mcoeff} lists; setting $M_{mn}=\ell^{2}\delta_{mn}$ reduces the same
expression to the isotropic Aifantis form~\cite{Aifantis1992}. Dropping the
memory and the flexoelectric tensor leaves the classical piezoelectric strip.
The present work differs from these in carrying the fractional memory on the
reduced modal coordinates and in treating the window's second moment, rather
than a scalar length, as the object that the discretization sees.

\FloatBarrier
